\chapter{Morfismos étales}
\label{I}

\marginpar{1}

Para simplificar la exposición, se supone que todos los preesquemas considerados son localmente noetherianos, al menos después del número~\Ref{I.2}.

\section{Nociones de cálculo diferencial}
\label{I.1}
Sea $X$ un preesquema sobre~$Y$, y sea $\Delta_{X/Y}$, o $\Delta$,
\label{indnot:ab}\oldindexnot{$\Delta_{X/Y}$ o simplemente $\Delta$|hyperpage}%
el morfismo diagonal $X\to X\times_Y X$. Es un morfismo de inmersión y, por tanto, un morfismo de inmersión cerrada de~$X$ en un abierto~$V$ de~$X\times_Y X$. Sea $\cal{I}_X$ el ideal del subpreesquema cerrado correspondiente a la diagonal en~$V$ (N.B.: si se quieren hacer las cosas de manera intrínseca, sin suponer que $X$ es separado sobre~$Y$ ---hipótesis que sería una tomadura de pelo---, habría que considerar la imagen inversa como haz de conjuntos de~$\cal{O}_{X\times X}$ sobre~$X$, y designar por $\cal{I}_X$ el ideal de aumentación de este últimoldots). El haz $\cal{I}_X/\cal{I}_X^2$ puede considerarse como un haz cuasicoherente sobre~$X$; se denota por~$\mathit{\Omega}^1_{X/Y}$.
\label{indnot:ac}\oldindexnot{$\mathit{\Omega}^1_{X/Y}$|hyperpage}%
Es de tipo finito si $X\to Y$ es de tipo finito. Se comporta bien respecto del cambio de base $Y'\to Y$. Se introducen asimismo los haces $\cal{O}_{X\times_Y X}/\cal{I}_X^{n+1}=\cal{P}^n_{X/Y}$,
\label{indnot:ad}\oldindexnot{$\cal{P}_{X/Y}^n$|hyperpage}%
que son haces de \emph{anillos} sobre~$X$ y dotan a~$X$ de una estructura de preesquema que puede denotarse por $\Delta_{X/Y}^n$
\label{indnot:ae}\oldindexnot{$\Delta_{X/Y}^n$|hyperpage}%
y llamarse el \emph{$n$-ésimo entorno infinitesimal de}~$X/Y$.
\index{entorno infinitesimal de~$X/Y$|hyperpage}%
El sorites correspondiente es de una trivialidad absoluta, aunque bastante largo\footnote{\cf EGA~IV~16.3.}; sería prudente no hablar de él hasta el momento en que pueda decirse algo útil al respecto, con los morfismos lisos.

